\chapter{Tưởng Mình Sẽ Nói Mà Không}
\chapterintro{Nhiều câu nói được chuẩn bị cả tuần nhưng tới đúng lúc thì lại tan rất nhanh}{Đây là kiểu dũng cảm hụt rất quen của tuổi học trò.}

\begin{lucbatbox}{Lục bát: Soạn sẵn trong đầu}
Suốt tuần em đã luyện rồi\
Gặp nhau sẽ nói một lời cho xong\
Đến khi người ấy quay trông\
Bao câu chuẩn bị như không có nhà\
Miệng cười, tim đập, mắt xa\
Rồi em đi mất như là chưa chi
\end{lucbatbox}

\begin{tutuyetbox}{Thất ngôn tứ tuyệt: Dũng cảm nửa đường}
Trong đầu thì rất hùng hồn\
Ra ngoài thực tế lại dồn lặng im\
Mới hay từ nghĩ sang tìm\
Một câu để nói có thêm cây cầu
\end{tutuyetbox}

\begin{ghichu}
Tưởng mình sẽ nói mà không là một kiểu hụt hẫng rất nhẹ ngoài mặt nhưng lại kéo dài trong đầu khá lâu.
\end{ghichu}

\ngancach

\begin{lucbatbox}{Lục bát: Lần sau chắc chắn}
Về nhà em tự dặn lòng\
Lần sau nhất định không còn run đâu\
Hôm sau gặp lại từ đầu\
Kịch bản cũ lại qua cầu y như\
Tuổi trẻ nhiều lúc chần chừ\
Chỉ một câu ngắn mà như leo đèo
\end{lucbatbox}

\begin{tutuyetbox}{Thất ngôn tứ tuyệt: Im lặng đi qua}
Người kia đi khuất cuối sân\
Em còn đứng lại với ngần ấy câu\
Nói ra thì nhẹ biết bao\
Mà sao lúc ấy nghẹn vào tận tim
\end{tutuyetbox}